\documentclass[10pt,a4paper]{amsart}
\usepackage[margin=19mm]{geometry}
\usepackage{amsmath,amssymb,mathtools}
\usepackage[scr=rsfs,cal=euler]{mathalfa}
\usepackage{xfrac}
\usepackage[unicode,hidelinks,bookmarks=false]{hyperref}
\newcommand{\ie}{i.e.\ }
\newcommand{\cf}{cf.\ }
\newcommand{\resp}{respectively\ }
\newcommand{\Subsection}{\subsection}
\providecommand{\nobreakdash}{\nobreak\hbox{-}}
\setlength{\emergencystretch}{2em}
\multlinegap0pt
\usepackage{amssymb}
\usepackage{enumerate}
\usepackage{mathtools}
\usepackage{xcolor}
\newtheorem{prop}{Proposition}
\newtheorem{lem}{Lemma}
\def\D{\mathbb D}
\title{Critical cyclicity in Dirichlet-type spaces on the bidisk}
\author{Athanasios Beslikas, Pouriya Torkinejad Ziarati}
\date{Source: arXiv:2609.09952v2 (CC BY 4.0)}
\begin{document}
\maketitle

\noindent\emph{Source and licence.} Critical cyclicity in Dirichlet-type spaces on the bidisk. Version arXiv:2609.09952v2 (CC BY 4.0), distributed by arXiv under the Creative Commons Attribution 4.0 International licence, http://creativecommons.org/licenses/by/4.0/, as recorded in the arXiv licence metadata for this exact version and shown on the abstract page https://arxiv.org/abs/2609.09952v2. Full licence notice: \texttt{licenses/arxiv-2609.09952-CC-BY-4.0.txt}.\par\smallskip

\noindent\textit{Editorial scope.} Counted unit: the article's lem lem:kernel-estimates (source0.tex lines 434--447). The article's complete original proof of it is delivered with it (source0.tex lines 448--515, one \texttt{proof} environment of 68 lines). No mathematics is written, repaired or re-derived: the statement and the proof are the article's own lines, in the article's own notation, and no retained line is substituted or re-flowed.  Retained uncounted as the source-local context the proof needs: lines 343--349; lines 404--407; lines 409--412.  Nothing was omitted from any retained span, so the package carries no omission record.  No retained line contains a citation, so the wrapper carries no bibliography and no bibliography entry.  No retained line contains a figure environment, an image-inclusion command or drawing code, so no image is delivered and the package declares no asset dependency.  Editorial formatting only, in addition: an AMS \texttt{amsart} preamble that restates the article's own declarations and macros used by the retained lines, namely the environments \texttt{prop}, \texttt{lem} on separate counters, together with the standard packages those lines need.  The retained environments print in this document as Proposition 1, Lemma 1 in order of appearance, whereas the article prints the same items with the numbering of its own full document.  The complete native source of the article as distributed in the arXiv e-print for this exact version is bundled unchanged as \texttt{sources/209863/source0.tex}.
\begin{prop} \label{prop: kernel}
Let $\beta \in (0,2)$. Then $\mathcal D_\beta(\mathbb{D}^2)$ is an RKHS (reproducing kernel Hilbert space), i.e., it can be endowed with an inner product for which the reproducing kernel $k_\beta$ is given by
\begin{equation*}
    k_{\beta}(z,w) = \int_0^1 \frac{dt}{\left(1 - t z_1 \overline{w}_1 - (1-t) z_2 \overline{w}_2 \right)^{2-\beta}}, 
    \quad z,w \in \mathbb{D}^2.
\end{equation*}
\end{prop}

\begin{equation}\label{eq:affine-power-estimate}
    \int_0^1|(1-t)x+t y|^{-\alpha}\,dt
    \lesssim_\alpha M^{-\alpha}.
\end{equation}

\begin{equation}\label{eq:affine-log-estimate}
    \int_0^1\log\frac{e}{|(1-t)x+ty|}\,dt
    \lesssim\log\frac{e}{M}.
\end{equation}

\begin{lem}\label{lem:kernel-estimates}
    Let $1<\beta\leq 2$. Then, for all $z,w\in \overline{\D}^2$, we have
    \begin{equation*}
        \operatorname{Re} k_\beta(z,w) \asymp |k_\beta(z,w)|.
    \end{equation*}
    Moreover, if $z,w\in \mathbb{T}^2$ and $1<\beta<2$, then
    \begin{equation*}
        |k_\beta(z,w)| \asymp \frac{1}{|z-w|^{2-\beta}}.
    \end{equation*}
    In the case $\beta=2$, we have
    \begin{equation*}
        |k_2(z,w)| \asymp \log \frac{e \sqrt{2}}{|z-w|}.
    \end{equation*}
\end{lem}

\begin{proof}
    The first part of the statement is an immediate consequence of Proposition~\ref{prop: kernel}. 

        We now assume that \(z,w\in\mathbb T^2\) and \(z\neq w\).
    Since the kernel is invariant under coordinate rotations (i.e., rotating each coordinate separately), we may
    assume that
    $w=(1,1)$, $z=(e^{i\theta},e^{i\varphi})$ and $\theta,\varphi\in[-\pi,\pi].$    When \(\theta\) and \(\varphi\) are sufficiently small,
    \begin{equation}
        |z-w|\asymp \max\{|\theta|,|\varphi|\}.
        \label{eq:angle-distance}
    \end{equation}

    Let \(1<\beta<2\) and put \(\alpha=2-\beta\). Then,
    \begin{equation}
        |k_\beta(z,w)|
        \asymp
        \int_0^1
        \left|
            t(1-e^{i\theta})+(1-t)(1-e^{i\varphi})
        \right|^{-\alpha}\,dt.
        \label{eq:kernel-power-integral}
    \end{equation}
    Uniformly in \(t\),
    $
        \left|
            t(1-e^{i\theta})+(1-t)(1-e^{i\varphi})
        \right|
        \lesssim\max\{|\theta|,|\varphi|\},
    $
    which gives the required lower bound in
    \eqref{eq:kernel-power-integral}. On the other hand,
    \[
        \left|
            t(1-e^{i\theta})+(1-t)(1-e^{i\varphi})
        \right|
        \geq
        |t\sin\theta+(1-t)\sin\varphi|.
    \]
    Since
    \[
        \max\{|\sin\theta|,|\sin\varphi|\}
        \asymp\max\{|\theta|,|\varphi|\},
    \]
    estimates \eqref{eq:affine-power-estimate} and
    \eqref{eq:angle-distance} give
    \[
        |k_\beta(z,w)|
        \asymp
        \frac{1}{|z-w|^{2-\beta}}.
    \]

    For \(\beta=2\), the same
    argument, using \eqref{eq:affine-log-estimate}, gives
    \[
    \begin{aligned}
        |k_2(z,w)|
        &\asymp
        \int_0^1
        \log\frac{e}{
        \left|t(1-e^{i\theta})+(1-t)(1-e^{i\varphi})\right|}
        \,dt \\
        &\asymp
        \log\frac{e}{\max\{|\theta|,|\varphi|\}}
        \asymp
        \log\frac{e\sqrt2}{|z-w|}.
    \end{aligned}
    \]
\end{proof}

\end{document}
